\documentclass[11pt,a4paper]{ctexart}

\usepackage[margin=2.45cm]{geometry}
\usepackage{amsmath,amssymb,mathtools}
\usepackage{booktabs,array,longtable}
\usepackage{enumitem}
\usepackage{xcolor}
\usepackage[hidelinks]{hyperref}

\setlength{\parindent}{2em}
\setlength{\parskip}{0.35em}
\setlist[enumerate]{itemsep=0.35em,topsep=0.4em}
\setlist[itemize]{itemsep=0.25em,topsep=0.35em}

\newcommand{\R}{\mathbb{R}}
\newcommand{\T}{\mathcal{T}}
\newcommand{\F}{\mathcal{F}}
\newcommand{\dif}{\mathop{}\!\mathrm{d}}

\title{第四题：Poisson 方程\\
有限体积空间离散与 OpenFOAM 数值验证\\
\large 自包含题目}
\author{}
\date{}

\begin{document}
\maketitle

\section{问题背景与符号约定}

本题严格采用原题第 4 题的方程符号
\begin{equation}
  \nabla\cdot\nabla\phi=\omega,
  \qquad\text{即}\qquad
  \Delta\phi=\omega.
  \label{eq:poisson-original-sign}
\end{equation}
参考文献~\cite{FerrerWillden2011} 使用的是
$-\Delta u+\alpha u=g$。因此，当 $\alpha=0$ 时，文献中的 Poisson
形式 $-\Delta u=g$ 与本题的关系为
\begin{equation}
  u=\phi,
  \qquad g=-\omega.
  \label{eq:sign-map}
\end{equation}
后文始终使用式~\eqref{eq:poisson-original-sign} 的符号，不混用两种约定。

Poisson 方程不含时间变量，所以原题所称的 ``semi-discretized form''
在本题中应理解为\textbf{有限体积空间离散后得到的代数方程组}，不存在时间离散。

\section{连续初边值问题}

令
\begin{equation}
  \Omega=(0,1)^2\subset\R^2,
  \qquad
  \partial\Omega=\Gamma_{x=0}\cup\Gamma_{x=1}
  \cup\Gamma_{y=0}\cup\Gamma_{y=1}.
\end{equation}
求标量场 $\phi:\overline\Omega\to\R$，使其满足
\begin{subequations}\label{eq:strong-problem}
\begin{align}
  \Delta\phi(x,y)&=\omega(x,y),
  &&(x,y)\in\Omega,\label{eq:strong-pde}\\
  \phi(x,y)&=g_D(x,y),
  &&(x,y)\in\partial\Omega.\label{eq:strong-bc}
\end{align}
\end{subequations}

采用制造解
\begin{equation}
  \phi_{\mathrm{exact}}(x,y)
  =\cos(\pi x)\cos(\pi y).
  \label{eq:exact}
\end{equation}
由式~\eqref{eq:exact} 直接计算得到
\begin{equation}
  \omega(x,y)
  =\Delta\phi_{\mathrm{exact}}(x,y)
  =-2\pi^2\cos(\pi x)\cos(\pi y).
  \label{eq:source}
\end{equation}
Dirichlet 边界数据必须取制造解在边界上的迹，即
\begin{equation}
  g_D=\phi_{\mathrm{exact}}\big|_{\partial\Omega}.
\end{equation}
为使问题完全自包含，四条边上的边界值明确写为
\begin{subequations}\label{eq:explicit-bc}
\begin{align}
  \phi(0,y)&= \cos(\pi y),
  &\phi(1,y)&=-\cos(\pi y),
  &&0\le y\le1,\\
  \phi(x,0)&= \cos(\pi x),
  &\phi(x,1)&=-\cos(\pi x),
  &&0\le x\le1.
\end{align}
\end{subequations}
四个角点的取值分别为
$\phi(0,0)=1$、$\phi(1,0)=-1$、$\phi(0,1)=-1$、
$\phi(1,1)=1$，故相邻边给出的 Dirichlet 数据在角点相容。

\section{网格、未知量与统一面记号}

分别构造两组逐级加密、与正方形边界相符的网格族：
\begin{equation}
  \{\T_h^{\mathrm{quad}}\}_{h\downarrow0},
  \qquad
  \{\T_h^{\mathrm{tri}}\}_{h\downarrow0},
\end{equation}
其中前者仅由四边形控制体组成，后者仅由三角形控制体组成。
每组至少使用三个逐级加密层级，并在报告中给出每级网格的单元数、
面数以及代表性网格尺度。不得把三角形和四边形网格上的误差混合后计算
收敛阶；两类网格应分别分析。

统一采用下列有限体积记号。

\begin{longtable}{>{\centering\arraybackslash}p{0.18\textwidth}p{0.72\textwidth}}
\toprule
记号 & 含义\\
\midrule
$c$ & 当前控制体；其中心为 $\mathbf{x}_c$，体积（二维单位厚度下为面积）为 $V_c=|\Omega_c|$。\\
$\F_c$ & 控制体 $c$ 的全部面集合。\\
$f$ & 一个内部面；$\mathbf{x}_f$ 为面心，$|S_f|$ 为面长度（二维单位厚度下的面面积）。\\
$O_f,N_f$ & 内部面 $f$ 的 owner 单元和 neighbour 单元。全局面方向固定为 $O_f\to N_f$。\\
$\mathbf n_f$ & 从 $O_f$ 指向 $N_f$ 的全局单位面法向量。\\
$\mathbf S_f$ & 全局有向面积向量，定义为 $\mathbf S_f=|S_f|\mathbf n_f$。\\
$\mathbf d_f$ & 中心连线向量 $\mathbf d_f=\mathbf x_{N_f}-\mathbf x_{O_f}$；$d_f=\|\mathbf d_f\|$。\\
$\sigma_{cf}$ & 面与单元的关联符号：当 $c=O_f$ 时为 $+1$，当 $c=N_f$ 时为 $-1$。\\
$\mathbf S_{cf}$ & 相对于单元 $c$ 的外向面积向量：$\mathbf S_{cf}=\sigma_{cf}\mathbf S_f=|S_f|\mathbf n_{cf}$。\\
$\phi_c$ & 单元中心 $\mathbf x_c$ 处的数值未知量，近似 $\phi(\mathbf x_c)$。\\
$\omega_c$ & 单元源项，取 $\omega_c=\omega(\mathbf x_c)$；若使用更高阶体积分，须明确说明。\\
$H_f$ & 沿全局方向 $O_f\to N_f$ 的内部面离散通量，近似 $\nabla\phi\cdot\mathbf S_f$。\\
$Q_c$ & 尚未除以 $V_c$ 的单元代数残差。\\
$R_c$ & 除以体积后的归一化残差，$R_c=Q_c/V_c$。\\
\bottomrule
\end{longtable}

边界面没有 neighbour 单元。对边界面 $f\subset\partial\Omega_c$，
$\mathbf n_{cf}$ 始终表示区域外法向量，
$\mathbf S_{cf}=|S_f|\mathbf n_{cf}$，并在面心施加
\begin{equation}
  \phi_f=g_D(\mathbf x_f).
  \label{eq:boundary-face-value}
\end{equation}

\section{有限体积空间离散要求}

\subsection{控制体积分形式}

对任意控制体 $\Omega_c$ 积分，并应用 Gauss 散度定理：
\begin{equation}
  \int_{\Omega_c}\Delta\phi\,\dif V
  =\int_{\partial\Omega_c}\nabla\phi\cdot\mathbf n_c\,\dif S
  =\int_{\Omega_c}\omega\,\dif V.
\end{equation}
采用单元中心源项积分后，有限体积平衡式为
\begin{equation}
  \sum_{f\in\F_c}
  (\nabla\phi)_f\cdot\mathbf S_{cf}
  =\omega_cV_c.
  \label{eq:fv-balance}
\end{equation}
要求从式~\eqref{eq:fv-balance} 出发推导内部面和 Dirichlet 边界面的
离散贡献，并说明同一个内部面通量只计算一次，随后按
\begin{equation}
  Q_{O_f}\mathrel{+}=H_f,
  \qquad
  Q_{N_f}\mathrel{-}=H_f
  \label{eq:owner-neighbour-assembly}
\end{equation}
装配，从而保证内部面通量在全局求和中严格抵消。

\subsection{OpenFOAM 基准空间格式}

为使计算设置唯一且可复现，本题的基准计算使用 OpenFOAM 的
\begin{center}
\texttt{laplacian(phi) Gauss linear corrected;}
\end{center}
即用 Gauss 定理离散 Laplacian 项，用线性插值处理面量，并对内部面
法向梯度加入显式非正交修正~\cite{OpenFOAMSchemes}。用于修正项的
单元梯度采用
\begin{center}
\texttt{grad(phi) Gauss linear;}
\end{center}

为明确该格式的数学对象，对内部面 $f$ 定义
\begin{equation}
  \mathbf e_f=\frac{\mathbf d_f}{d_f},
  \qquad
  \alpha_f=\frac{1}{\mathbf n_f\cdot\mathbf e_f},
  \qquad
  \mathbf k_f=\mathbf n_f-\alpha_f\mathbf e_f.
  \label{eq:nonorth-decomp}
\end{equation}
其中要求 $\mathbf n_f\cdot\mathbf e_f>0$。则 corrected
面法向梯度可写为
\begin{equation}
  (\nabla\phi)_f\cdot\mathbf n_f
  \approx
  \alpha_f\frac{\phi_{N_f}-\phi_{O_f}}{d_f}
  +(\nabla\phi)_f^{\mathrm{lin}}\cdot\mathbf k_f,
  \label{eq:corrected-sngrad}
\end{equation}
从而全局内部面通量为
\begin{equation}
  H_f
  =|S_f|\left[
  \alpha_f\frac{\phi_{N_f}-\phi_{O_f}}{d_f}
  +(\nabla\phi)_f^{\mathrm{lin}}\cdot\mathbf k_f
  \right].
  \label{eq:global-face-flux}
\end{equation}
当网格正交时，$\mathbf n_f=\mathbf e_f$，故
$\alpha_f=1$、$\mathbf k_f=\mathbf0$，式~\eqref{eq:global-face-flux}
退化为标准中心差分通量
\begin{equation}
  H_f=\frac{|S_f|}{d_f}
  (\phi_{N_f}-\phi_{O_f}).
\end{equation}

Gauss linear 单元梯度应从
\begin{equation}
  (\nabla\phi)_c
  \approx\frac{1}{V_c}
  \sum_{f\in\F_c}\phi_f^{\mathrm{lin}}\mathbf S_{cf}
  \label{eq:gauss-gradient}
\end{equation}
构造；内部面梯度 $(\nabla\phi)_f^{\mathrm{lin}}$ 再由相邻单元梯度
线性插值得到。边界面使用式~\eqref{eq:boundary-face-value} 给出的
精确 Dirichlet 面值，并采用 OpenFOAM 的 \texttt{fixedValue}
边界离散。实现时可把各边面心处的 $g_D(\mathbf x_f)$ 写为
\texttt{nonuniform List<scalar>}，或使用等价的解析边界实现；不得把
四条边错误地设为同一个常数。

\subsection{单元残差与线性系统}

对控制体 $c$，定义未归一化残差
\begin{equation}
  Q_c(\phi)
  =\sum_{f\in\F_c^{\mathrm{int}}}\sigma_{cf}H_f
  +\sum_{f\in\F_c^{\mathrm{bnd}}}H_{cf}^{\mathrm{bnd}}
  -\omega_cV_c,
  \label{eq:Qc}
\end{equation}
以及归一化残差
\begin{equation}
  R_c(\phi)=\frac{Q_c(\phi)}{V_c}.
  \label{eq:Rc}
\end{equation}
离散问题是求所有单元未知量 $\{\phi_c\}$，使得
\begin{equation}
  Q_c(\phi)=0
  \quad\Longleftrightarrow\quad
  R_c(\phi)=0,
  \qquad \forall c\in\T_h.
  \label{eq:algebraic-problem}
\end{equation}
要求将其装配为线性系统
\begin{equation}
  A\boldsymbol\phi=\mathbf b
\end{equation}
并求解。为避免正负号混淆，OpenFOAM 中可以直接实现
\begin{equation}
  -\Delta\phi=-\omega
  \label{eq:spd-equivalent}
\end{equation}
这一与原题完全等价的形式，即使用
\begin{center}
\texttt{-fvm::laplacian(phi) == -omega}
\end{center}
构造具有通常正定符号约定的线性系统。必须在报告中说明采用的是
式~\eqref{eq:poisson-original-sign} 还是等价式~\eqref{eq:spd-equivalent}，
并验证源项符号与之相容。

\section{数值实验与评价指标}

\subsection{必须完成的计算}

分别在四边形网格族 $\T_h^{\mathrm{quad}}$ 和三角形网格族
$\T_h^{\mathrm{tri}}$ 上完成以下任务：

\begin{enumerate}[label=\arabic*.]
  \item 计算离散解 $\phi_h$，绘制 $\phi_h$ 的等值线图或云图，并与
  $\phi_{\mathrm{exact}}$ 使用相同的色标进行比较。

  \item 绘制点误差场
  \begin{equation}
    e_c=\phi_c-\phi_{\mathrm{exact}}(\mathbf x_c).
  \end{equation}

  \item 对每个网格层级报告单元数 $N_e$、代表性尺度
  \begin{equation}
    h=N_e^{-1/2},
    \label{eq:effective-h}
  \end{equation}
  线性求解器及收敛容差，并计算下述 $L_1$、$L_2$ 和 $L_\infty$
  相对误差。

  \item 分别对三角形网格和四边形网格计算相邻层级之间的观测收敛阶，
  作 $\log E$--$\log h$ 图，并比较两类网格的误差常数与收敛行为。

  \item 检查最终代数残差，并说明若收敛阶偏离预期，是由网格非正交性、
  边界离散、源项积分、线性求解误差还是实现错误导致。
\end{enumerate}

\subsection{误差范数}

令 $N_e$ 为单元总数，
\begin{equation}
  \phi_c^{n}=\phi_c,
  \qquad
  \phi_c^{e}=\phi_{\mathrm{exact}}(\mathbf x_c),
\end{equation}
其中上标 $n$ 表示 numerical，而不是时间层。严格按原题附录 A 的
体积加权相对误差定义
\begin{align}
  E(L_1)
  &=\frac{\displaystyle\sum_{c=1}^{N_e}
  |\phi_c^{n}-\phi_c^{e}|V_c}
  {\displaystyle\sum_{c=1}^{N_e}|\phi_c^{e}|V_c},
  \label{eq:L1}\\[0.5em]
  E(L_2)
  &=\left(
  \frac{\displaystyle\sum_{c=1}^{N_e}
  (\phi_c^{n}-\phi_c^{e})^2V_c}
  {\displaystyle\sum_{c=1}^{N_e}(\phi_c^{e})^2V_c}
  \right)^{1/2},
  \label{eq:L2}\\[0.5em]
  E(L_\infty)
  &=\frac{\displaystyle\max_{1\le c\le N_e}
  |\phi_c^{n}-\phi_c^{e}|}
  {\displaystyle\max_{1\le c\le N_e}|\phi_c^{e}|}.
  \label{eq:Linf}
\end{align}

对二维网格上单元数为 $N_1<N_2$ 的两个相邻层级，误差 $E_q$
（$q\in\{1,2,\infty\}$）的观测收敛阶定义为
\begin{equation}
  p_q
  =\frac{\log\!\left(E_q(N_1)/E_q(N_2)\right)}
  {\log\!\left((N_2/N_1)^{1/2}\right)}
  =\frac{\log\!\left(E_q(h_1)/E_q(h_2)\right)}
  {\log(h_1/h_2)}.
  \label{eq:order}
\end{equation}

\section{最终交付要求}

最终提交应至少包含：
\begin{enumerate}[label=(\arabic*)]
  \item 从强形式到式~\eqref{eq:fv-balance} 的有限体积积分推导；
  \item 式~\eqref{eq:global-face-flux} 的正交项与非正交修正项，以及
  owner--neighbour 的守恒装配；
  \item Dirichlet 边界数据、源项和方程符号的一致性说明；
  \item OpenFOAM 求解器或可复现算例，包含 \texttt{fvSchemes}、
  \texttt{fvSolution}、网格和场文件；
  \item 三角形与四边形逐级加密网格上的解场、误差场、
  $E(L_1)$、$E(L_2)$、$E(L_\infty)$ 及观测收敛阶；
  \item 对守恒性、代数残差、网格影响和异常收敛行为的解释。
\end{enumerate}

\paragraph{参考资料使用边界。}
参考文献~\cite{FerrerWillden2011} 采用高阶 DG--SIPG 方法；本题只从中
继承制造解、区域、Dirichlet 数据及 Poisson/Helmholtz 验证问题的设定，
不把该文的 DG 罚参数或 $p$ 阶理论收敛率移植到 OpenFOAM 的单元中心
有限体积格式中。有限体积离散必须以本题规定的 OpenFOAM 基准空间格式
为准。

\begin{thebibliography}{9}

\bibitem{OriginalTests}
\emph{Numerical TESTs for Partial Differential Equations (PDEs)},
第四题 ``Poisson equation'' 与附录 A，2023-09-26（用户提供文件）。

\bibitem{FerrerWillden2011}
E. Ferrer and R. H. J. Willden,
``A high order discontinuous Galerkin finite element solver for the
incompressible Navier--Stokes equations,''
\emph{Computers \& Fluids}, 46(1):224--230, 2011.
\href{https://doi.org/10.1016/j.compfluid.2010.10.018}
{doi:10.1016/j.compfluid.2010.10.018}.

\bibitem{OpenFOAMSchemes}
OpenFOAM Foundation,
``OpenFOAM User Guide: Numerical schemes,''
\url{https://doc.cfd.direct/openfoam/user-guide/fvschemes}.

\end{thebibliography}

\end{document}
